% Owner: TV3 (P5.3b) -- Reviewer: TV1
% Tables T2, T6, T7, and B4 are generated by `python -m experiments.report_tables`.
\subsection{Alpha-Beta on Gomoku, Rollout Policies, and Head-to-Head Play}
\label{subsec:results_tv3}

\textbf{Ablation of alpha-beta on Gomoku (RQ2).} Table~\ref{tab:t2} switches off one component at a time and plays the reduced agent against the full agent on the $15\times 15$ board, with one second per move and $50$ games per pair. Move ordering is decisive: without it the agent loses all $50$ games. Since the search keeps only $16$ moves per node, an unsorted move list makes this cut keep arbitrary moves, so the agent searches more nodes ($2\,794$ per move against $583$) but the wrong ones. Iterative deepening is the next most important component ($8\%$ wins): without it, a search that runs out of time has no completed depth to fall back on. Searching all empty cells instead of the candidates near the stones also loses $44$ of $50$ games, because the branching factor rises from at most $16$ to over $200$. Switching off the immediate win/block test or the transposition table gives $48\%$ and $44\%$, and neither difference is significant. A likely explanation is that the pattern ordering already ranks winning and blocking moves first, and that few positions repeat at the depths reached within one second.

\begin{table}[t]
\centering
\caption{T2: Alpha-beta with one component switched off (A) against the full agent on $15\times 15$ Gomoku, $1$~s per move, $50$ games per row}
\label{tab:t2}
\resizebox{\columnwidth}{!}{% Generated by experiments/report_tables.py -- do not edit by hand.
\begin{tabular}{lrrrr}
\toprule
Variant (A) & W/D/L & Win \% [95\% CI] & Nodes/move A & Nodes/move full \\
\midrule
No candidate restriction & 6/0/44 & 12.0 [5.6, 23.8] & 524 & 539 \\
No win/block detection & 24/0/26 & 48.0 [34.8, 61.5] & 899 & 572 \\
No pattern move ordering & 0/0/50 & 0.0 [0.0, 7.1] & 2\,794 & 583 \\
No iterative deepening & 4/4/42 & 8.0 [3.2, 18.8] & 673 & 656 \\
No transposition table & 22/0/28 & 44.0 [31.2, 57.7] & 577 & 542 \\
\bottomrule
\end{tabular}
}
\end{table}

\textbf{Rollout policies (RQ6).} Table~\ref{tab:t6} compares the rollout policies on $9\times 9$ Gomoku. Knowledge makes each simulation more expensive: L1 and L2 run at about half the speed of L0, the softmax sampling of L3 at about one twelfth, and L4 is twice as fast as L3 because it stops after eight moves. With an equal number of iterations ($200$), every policy beats L0 ($p < 0.001$), and the gain grows with knowledge, from $65.5\%$ for L2 to $90\%$ for L4. With equal time per move the order changes: L1, L2, and L3 now lose to L0 or tie with it, because L0 runs several times more simulations in the same time, and only L4, which combines knowledge with a cheap cut-off, still wins ($67\%$ at $0.08$~s per move, $p < 0.001$). The same pattern appears on trap positions. On $3\times 3$ trap E1, L1--L4 block the threat with only $10$ iterations in all $200$ runs, while L0 needs $50$. On a $9\times 9$ position where the opponent has a four, L1--L4 block it in every run from $60$ iterations, while L0 finds the block in at most $29\%$ of the runs with $100$ iterations. Against an open three, which is a threat two moves deep, the tactical rule of L1 and L2 no longer applies: with $120$ iterations the correct block is found in $12\%$ (L0), $26\%$ (L2), $48\%$ (L3), and $61\%$ (L4) of the runs.

\begin{table}[t]
\centering
\caption{T6: Rollouts per second, and win rate (\%) of MCTS with each rollout policy against MCTS with L0 on $9\times 9$ Gomoku at equal iterations or equal time, $200$ games per cell}
\label{tab:t6}
\resizebox{\columnwidth}{!}{% Generated by experiments/report_tables.py -- do not edit by hand.
\begin{tabular}{lrrrrr}
\toprule
Rollout & Sims/s $9\times 9$ & Sims/s $15\times 15$ & 200 it. & 0.08 s & 0.5 s \\
\midrule
L0 & 29\,502 & 13\,208 & -- & -- & -- \\
L1 & 16\,936 & 8\,698 & 75.0 & 41.5 & 48.0 \\
L2 & 14\,222 & 8\,334 & 65.5 & 44.0 & 43.0 \\
L3 & 2\,455 & 2\,434 & 85.5 & 35.0 & 29.0 \\
L4 & 4\,843 & 4\,438 & 90.0 & 67.0 & 55.0 \\
\bottomrule
\end{tabular}
}
\end{table}

\textbf{Alpha-beta versus MCTS on $15\times 15$ (RQ6).} Table~\ref{tab:t7} plays alpha-beta against MCTS with L2 or L4 rollouts ($c = 0.5$, tree reuse) at $0.5$, $1$, and $2$ seconds per move, $100$ games per cell with the first move alternated. Alpha-beta wins $594$ of the $600$ games. All six MCTS wins come when MCTS moves first, and three of them with L4 at $2$ seconds per move. MCTS adds more nodes per move than alpha-beta visits, but its random simulations on a board of $225$ cells are too noisy to recognize the threats that the pattern evaluation of alpha-beta scores directly.

\begin{table}[t]
\centering
\caption{T7: Alpha-beta (AB) against MCTS on $15\times 15$ Gomoku at equal time per move; AB wins out of $50$ games when each side moves first, and nodes per move}
\label{tab:t7}
\resizebox{\columnwidth}{!}{% Generated by experiments/report_tables.py -- do not edit by hand.
\begin{tabular}{rlrrrrr}
\toprule
s/move & MCTS & AB W/D/L & AB first & MCTS first & AB nodes & MCTS nodes \\
\midrule
0.5 & L2 & 100/0/0 & 50/50 & 50/50 & 417 & 1\,228 \\
0.5 & L4 & 99/0/1 & 50/50 & 49/50 & 394 & 690 \\
1 & L2 & 99/0/1 & 50/50 & 49/50 & 681 & 2\,247 \\
1 & L4 & 99/0/1 & 50/50 & 49/50 & 636 & 1\,271 \\
2 & L2 & 100/0/0 & 50/50 & 50/50 & 1\,241 & 4\,417 \\
2 & L4 & 97/0/3 & 50/50 & 47/50 & 1\,262 & 2\,313 \\
\bottomrule
\end{tabular}
}
\end{table}

\textbf{Simulation improvements (bonus B4).} Table~\ref{tab:b4} measures the two improvements of Section~\ref{subsec:rollouts} on $9\times 9$ Gomoku, $200$ games per row. The extended forced moves make every rollout policy stronger at equal iterations ($65$--$85\%$ wins against the same policy without them) and also at equal time ($71$--$84.5\%$), although they slow the L2 and L3 rollouts by about a factor of three. On the open-three trap position they raise the rate at which MCTS finds the correct block from $10\%$ to $72.5\%$ for L2 and from $57.5\%$ to $95\%$ for L4. Root parallelization with $K=2$ or $K=4$ trees beats a single tree at the same wall-clock time ($59\%$ and $61.5\%$), and with $120$ iterations per tree ($62.5\%$ and $71.5\%$); the gain at equal time is modest because the independent trees share no statistics below the root.

\begin{table}[t]
\centering
\caption{B4: Win rate of MCTS with the extended forced moves (L$x$-B4) against the same policy without them, and of root parallelization with $K$ trees against one tree, on $9\times 9$ Gomoku}
\label{tab:b4}
\resizebox{\columnwidth}{!}{% Generated by experiments/report_tables.py -- do not edit by hand.
\begin{tabular}{llr}
\toprule
Comparison & Budget & Win \% [95\% CI] \\
\midrule
L2-B4 vs L2 & 200 iterations & 85.0 [79.4, 89.3] \\
L3-B4 vs L3 & 200 iterations & 74.5 [68.0, 80.0] \\
L4-B4 vs L4 & 200 iterations & 65.0 [58.2, 71.3] \\
L2-B4 vs L2 & 0.08 s/move & 74.5 [68.0, 80.0] \\
L3-B4 vs L3 & 0.08 s/move & 84.5 [78.8, 88.9] \\
L4-B4 vs L4 & 0.08 s/move & 71.0 [64.4, 76.8] \\
$K=2$ vs $K=1$ & 0.04 s/move & 59.0 [52.1, 65.6] \\
$K=2$ vs $K=1$ & 120 it./worker & 62.5 [55.6, 68.9] \\
$K=4$ vs $K=1$ & 0.04 s/move & 61.5 [54.6, 68.0] \\
$K=4$ vs $K=1$ & 120 it./worker & 71.5 [64.9, 77.3] \\
\bottomrule
\end{tabular}
}
\end{table}

\textbf{Xiangqi.} On Xiangqi, alpha-beta with a material and piece-square evaluation and MCTS with an L4-style rollout both beat a random player in all $60$ games. Alpha-beta beats MCTS in all $60$ games at both $0.1$ and $0.2$ seconds per move. At $0.1$ seconds MCTS completes only about $100$ simulations per move, too few for a game with around $40$ legal moves per position.
